\section{Viabilidade Corporativa e o Ponto de Operação Industrial (QP4)}
\label{sec:viabilidade_corporativa}

Somente a qualidade preditiva, não garante a viabilidade operacional de um modelo em um ambiente agroindustrial. A resposta à Questão de Pesquisa QP4, estruturada empiricamente na Seção \ref{sec:viabilidade_industrial}, atesta que o verdadeiro estado da arte para a tipagem do algodão não reside, obrigatoriamente, na arquitetura com o maior \textit{Matthews Correlation Coefficient} (MCC) bruto, mas sim naquela que otimiza o \textit{trade-off} entre assertividade diagnóstica e restrições de \textit{hardware}, avaliada por meio de uma análise comparativa/descritiva de ranking de eficiência industrial, sem hipótese nula associada.

A análise topológica demonstrou que a DenseNet, embora seja a líder matemática em capacidade de generalização espacial, impõe uma penalidade de latência. O seu padrão de conectividade densa, onde cada camada recebe mapas de características de todas as camadas precedentes, gera um alto custo de acesso à memória.. No contexto do chão de fábrica, onde a pluma de algodão é processada em esteiras de fluxo contínuo de alta velocidade, uma latência de inferência de $19,04$ ms por \textit{patch} visual pode inviabilizar a inspeção volumétrica em tempo real ou, alternativamente, exige um investimento financeiro proibitivo em \textit{clusters} de processamento gráfico (GPUs).

Em contrapartida, a consagração da ConvNeXt como o ``ponto de operação'' (\textit{sweet spot}) ilustra a maturidade do paradigma convolucional moderno. A rede demonstrou ser capaz de sustentar o mesmo rigor estatístico de classificação da líder (evidenciado pelo desempenho estatisticamente semelhante no teste de McNemar) operando em apenas $6,44$ ms. Esta agilidade é fruto de refinamentos de macro e micro-design, como o uso estratégico de convoluções \textit{depthwise} com \textit{kernels} amplos ($7 \times 7$), que simulam a amplitude do campo receptivo dos \textit{Transformers}, mas preservando a eficiência computacional das CNNs.

A identificação empírica deste balanço transcende a contribuição algorítmica puramente abstrata e adentra o domínio prático da Engenharia de Sistemas. Ao eleger a ConvNeXt sob iluminação otimizada (RGB/BR) como a solução definitiva, esta pesquisa transpõe a barreira teórica e entrega à indústria uma ferramenta de \textit{Edge AI} madura. O sistema final atende simultaneamente às rigorosas exigências morfológicas da Instrução Normativa nº 24/2016 e aos requisitos de escalabilidade, cadência e eficiência energética mandatórios para a automação sustentável do agronegócio. Embora as métricas de latência tenham sido aferidas em arquitetura Apple Silicon (M3), estima-se que a ConvNeXt mantenha um \textit{throughput} compatível em \textit{hardwares} industriais padrão.